% Appendix: observations made while building the systems that bear on the reported results.
\section{Development Observations}
\label{app:observations}

% The stored prompts are from a commercial development series: no value from it is printed (README, "Rules
% that are easy to break"; Ethical Considerations). The OBS-bake-* keys stay declared and unused.
\paragraph{Choice of generator.} We sent stored claim-extraction prompts from a commercial
development series to three Gemini Flash models, then checked the outputs of two of them,
3.1 Flash-Lite and 3.6 Flash, against their cited paragraphs. No numerical results from that series are reported
(Appendix~\ref{app:corpus}). Qualitatively, Gemini~3.1 Flash-Lite returned more facts than
Gemini~3.6 Flash, with a lower share fully supported and some unsupported claims; no
unsupported claims were observed in the reviewed 3.6 Flash outputs. We selected 3.6 Flash
for all stages and comparison systems. Model choice also affects entity discovery: from
the same volume-one Anne candidate list, Flash-Lite retained \res{OBS-gaz-lite} entities
and 3.6 Flash retained \res{OBS-gaz-f36}. This illustrates an effect on the entity inventory
before page generation.

\paragraph{Stories within the story.} Characters tell stories, invent companions and make
false confessions. Entity discovery admitted figures from embedded stories, including a
story-club heroine and a child's imagined cave lady, and verification accepted claims about
them (Appendix~\ref{app:evaluator}). Local evidence can therefore be insufficient to
distinguish the main narrative from an embedded or imagined account.

\paragraph{Reveal volumes need an earliest-statement check.} Disclosure assessment found
\res{OBS-gold-late} draft facts assigned reveal volumes later than their earliest supporting
statements. Pages containing these facts were consequently eligible. Reveal-volume
annotation requires searching for the earliest support, rather than retaining the first
passage encountered by the annotator.

\paragraph{Implicit temporal progression.}
Gilbert's age, \emph{nearly fourteen}, is extracted from volume one and remains on his
cutoff-4 page alongside the later marriage (Figure~\ref{fig:qual-overview}c). No extracted
age supersedes it, so its interval stays open. The graph records explicit changes; it does
not advance ages with narrative time. Reveal volumes order evidence rather than elapsed
story years. A historical age should therefore be marked as an observation at a particular
story time unless a supported update is available.

\paragraph{Cross-claim inference.}
Some updates require combining claims rather than extracting one explicit statement.
As an illustrative example, suppose Anne is two years younger than Gilbert and marries
him at age 25. The premises imply that Gilbert is 27 at their marriage. These ages are
hypothetical, not verified annotations of the displayed example. The current pipeline
does not systematically derive such updates. An extension would need to retain the
supporting premises, align their narrative times and expose a derived claim only after
all required premises become available at the reader's cutoff.

\paragraph{Coreference across extraction windows.}
Pronouns and other referring expressions can link a character or fact across sentences
and scenes without repeating a name. Extraction may miss these links or attribute a
claim to the wrong person when the relevant context lies outside its window. Both backward
and forward references require discourse context, while a later identity revelation must
remain unavailable at an earlier cutoff. These are qualitative failure modes; no error
rate for them is measured here.
